\DocumentMetadata{
  pdfversion=2.0,
  pdfstandard=ua-2,
  lang=en-US,
  tagging=on,
  tagging-setup={math/setup=mathml-SE,tabsorder=structure}
}
\documentclass[11pt,letterpaper]{article}
\usepackage[margin=0.68in,includefoot]{geometry}
\usepackage{newcomputermodern}
\usepackage{amsmath,amscd,mathtools}
\usepackage{tikz,adjustbox}
\providecommand{\square}{\mdlgwhtsquare}
\usepackage[hidelinks]{hyperref}
\hypersetup{
  pdftitle={Analysis PhD Exam, January 1988},
  pdfauthor={University of Florida Department of Mathematics},
  pdfsubject={Graduate mathematics examination},
  pdfdisplaydoctitle=true,
  bookmarksopen=true
}
\setcounter{secnumdepth}{0}
\setlength{\parindent}{0pt}
\setlength{\parskip}{0.28em}
\setlength{\emergencystretch}{3em}
\setlength{\leftmargini}{2.15em}
\setlength{\leftmarginii}{2.65em}
\setlength{\leftmarginiii}{3em}
\renewcommand{\labelenumi}{\textbf{\theenumi.}}
\pagestyle{empty}
\raggedbottom
\allowdisplaybreaks
\newenvironment{examproblems}[1]{%
  \begin{enumerate}%
  \setcounter{enumi}{#1}%
  \setlength{\topsep}{0.35em}%
  \setlength{\partopsep}{0pt}%
  \setlength{\itemsep}{0.48em}%
  \setlength{\parsep}{0.18em}%
}{\end{enumerate}}
\newenvironment{examsubparts}{%
  \begin{enumerate}%
  \setlength{\topsep}{0.18em}%
  \setlength{\partopsep}{0pt}%
  \setlength{\itemsep}{0.16em}%
  \setlength{\parsep}{0.1em}%
}{\end{enumerate}}
\newenvironment{examinstructions}{%
  \par\begingroup\itshape
}{%
  \par\endgroup\vspace{0.18em}
}
\makeatletter
\renewcommand\section{\@startsection{section}{1}{0pt}%
  {-1.25ex plus -.25ex minus -.1ex}%
  {0.55ex plus .1ex}%
  {\normalfont\large\bfseries}}
\renewcommand{\@maketitle}{%
  \newpage\null\vskip -1em%
  {\footnotesize\bfseries
    \href{https://www.ufl.edu/}{University of Florida}, \href{https://math.ufl.edu/}{Department of Mathematics}\par}%
  \vskip -0.5em%
  \tagstructbegin{tag=Title}%
    {\LARGE\bfseries\href{https://gma.math.ufl.edu/exams/analysis-phd/}{Analysis PhD Exam}, January 1988\par}%
  \tagstructend%
  \vskip 0.25em%
  \tagmcbegin{artifact}%
    \hrule height 1pt%
  \tagmcend%
  \vskip 1em%
}
\makeatother
\title{Analysis PhD Exam, January 1988}
\author{}
\date{}
\begin{document}
\maketitle
\thispagestyle{empty}
\begin{examinstructions}
Answer all six questions. Justify all work. To obtain any partial credit, all work must be presented in a neat and logical fashion.
\end{examinstructions}
\begin{examproblems}{0}
\item
State and prove the Radon–Nikodym Theorem.
\item
State and prove the Hahn–Banach Theorem.
\item
Let~\(\Omega\) be a topological space, \(\mathcal B(\Omega)\) the Borel subsets of~\(\Omega\). Suppose that for each~\(n\), \(\mu_n\) is a countably additive positive finite regular measure defined on \(\mathcal B(\Omega)\). Let 
\[
\mu(\mathord\cdot)=\sum_{n=1}^{\infty}\frac{\mu_n(\mathord\cdot)}{2^n\bigl(1+\mu_n(\Omega)\bigr)}.
\]
 Show that~\(\mu\) is also a regular countably additive measure.
\item
Let \(\{f_n:n\geq 0\}\) and \(\{g_n:n\geq -1\}\) be two sequences of real-valued measurable functions on a measurable space \((\Omega,\mathcal F)\). Let \(B\subset\mathbb R^1\) be a Borel set on \(\mathbb R^1\). Define 
\[
T(w)=
\begin{cases}
\sup\{n\geq 0:f_n(w)\in B\},&\text{if }\{n\geq 0:f_n(w)\in B\}\text{ is bounded and nonempty},\\
0,&\text{if }\{n\geq 0:f_n(w)\in B\}=\varnothing,\\
-1,&\text{if }\{n\geq 0:f_n(w)\in B\}\text{ is unbounded}.
\end{cases}
\]
\begin{examsubparts}
\item[\textnormal{(A)}]
Prove~\(T\) is \(\mathcal F\)-measurable.
\item[\textnormal{(B)}]
Define \(h(w)=g_{T(w)}(w)\); prove~\(h\) is \(\mathcal F\)-measurable.
\end{examsubparts}
\item
Let \(T:(-\pi/2,\pi/2)\to\mathbb R\) be given by \(T(x)=\tan x\). Let~\(\lambda\) be Lebesgue measure on \((-\pi/2,\pi/2)\) and define a measure~\(\mu\) on \((\mathbb R,\mathcal B(\mathbb R))\) by setting \(\mu(A)=\lambda(T^{-1}(A))\). Compute \(\dfrac{d\mu}{dx}\).
\item
Fix \(1\leq p<\infty\), and let \(f\in L^p(X,\mathcal A,\mu)\). Show that 
\[
\lVert f\rVert_p=\sup\left\{\int fg\,d\mu:\lVert g\rVert_q\leq 1\right\}.
\]
 Hint:
\begin{examsubparts}
\item[\textnormal{1.}]
Show that if \(\lVert g\rVert_q\leq 1\), then \(\int fg\,d\mu\leq\lVert f\rVert_p\).
\item[\textnormal{2.}]
Show that \(g=|f|^{p-1}\in L^q\).
\item[\textnormal{3.}]
State why doing 1. and 2. solves the problem.
\end{examsubparts}
\end{examproblems}
\end{document}
